% ===================================================================
%  TABELUL Nr. 8 — Secerișul. — Vânătoarea, Culesul viei, Pescuitul.
%  Traduction 2026 d'apres texto/io, controlee sur texto/fr, texto/en
%  et texto/es. Consigne : voir texto/ro/10-tabelul-01.tex.
%
%  L'appel de note est dans le TITRE de la scene — « Secerișul (*) ».
%  LE BLOC c2-03-1 EST UN ECART DECLARE dans renvoji.py.
%
%  LE TABLEAU DES MOISSONS PROLONGE CELUI DE LA FERME : « secera » la
%  faucille, « grâu » le ble, « spic » l'epi, « snop » la gerbe,
%  « arie » l'aire — les quatre premiers latins, le dernier grec par le
%  latin. La ou le roumain quitte le latin, c'est pour l'outil moderne,
%  jamais pour le geste ancien.
% ===================================================================

%%K t08-noto-1 noto
\VUnotes{143.2mm}{%
\textit{(*) Fiindcă acest cuvânt nu este întocmai potrivit: secerișul inului,
culesul viei, culesul merelor? Poate ar fi fost mai bine să luăm «}
\VUgras{mesono} \textit{», propus în} Mondo \textit{XIV, p. 242. Dar deocamdată
această rădăcină nouă nu a fost primită de Academie și așteaptă dezbaterea după
regulile obișnuite ale mișcării ido.}%
}

%%K t08-apar-1 apar
\VUsaut{5.0mm}
\VUtitre{15.0pt}{60}{TABELUL Nr. 8}
\VUsaut{3.0mm}
\VUcentre{13.2pt}{90}{\textit{Prima scenă.}}
\VUsaut{3.0mm}
\VUpk{10.2pt}{40}{Secerișul (*).}
\VUsaut{4.0mm}
\VUpk{10.2pt}{40}{Chipul câmpurilor.}
\VUsaut{3.0mm}

%%K t08-c1-01-1 p
1. --- Odată cu venirea \VUgras{verii} câmpurile și-au schimbat din nou chipul.
Sămânța încredințată brazdei a răsărit; din pământ s-a ridicat o plantă verde mică
și și-a dat spicul. Bobul s-a umplut, apoi s-a copt, și acum nu mai rămâne decât
să se strângă secerișul.

%%K t08-c1-02-1 p
2. --- \VUgras{Florile de câmp} s-au deschis. \VUgras{Albăstrelele}
\textsuperscript{(1)}, \VUgras{macii} \textsuperscript{(2)} și
\VUgras{degețelul} \textsuperscript{(3)} aduc în tonul neted al grânelor o ceartă
veselă cu \VUgras{corolele} lor proaspete.

%%K t08-c1-03-1 p
3. --- Pomii s-au îmbrăcat în frunze; roadele s-au copt. Micul \VUgras{cireș}
\textsuperscript{(4)} se pleacă sub \VUgras{cireșele} frumoase \textsuperscript{(5)},
pe care copiii le rup și le mănâncă cu desfătare.

%%K t08-c1-04-1 p
4. --- Turma poate fi acum toată ziua în voie.

%%K t08-c1-04-2 p
\VUgras{Mieii} \textsuperscript{(6)} au crescut și au ajuns \VUgras{oi}
frumoase \textsuperscript{(7)} și \VUgras{berbeci} frumoși \textsuperscript{(8)}
cu lână deasă. Pasc liniștiți \VUgras{iarba} pe \VUgras{pășune}
\textsuperscript{(9)}, presărată cu \VUgras{margarete}.

%%K t08-c1-04-3 p
Curând acestor animale li se va tunde \VUgras{lâna}, din care se face postavul
pentru hainele noastre.

%%K t08-tit-2 sub
\VUsaut{5.0mm}
\VUpk{10.2pt}{40}{Păstorul.}
\VUsaut{3.4mm}

%%K t08-c1-05-1 p
5. --- \VUgras{Păstorul} \textsuperscript{(10)}, pare-se, nu se uită prea mult la
ele. Îl privește numai furtuna care trece la zare, iar \VUgras{ciobănașul}
\textsuperscript{(13)}, ducându-și \VUgras{plosca} la gură \textsuperscript{(14)},
pare și mai fără grijă.

%%K t08-c1-06-1 p
6. --- Zăpușeala apasă; iată și \VUgras{câinele} \textsuperscript{(15)} care merge
să se răcorească; lipăie apa proaspătă a \VUgras{pârâului} \textsuperscript{(16)}
și sperie o biată \VUgras{broască} mică \textsuperscript{(17)}, pitită în iarba de
pe mal. Aceasta sare în apa limpede, unde înfloresc \VUgras{nuferii}
\textsuperscript{(18)}.

%%K t08-c1-07-1 p
7. --- \VUgras{Codobatura} \textsuperscript{(19)} se scaldă lângă micul
\VUgras{izvor} \textsuperscript{(20)}, care țâșnește între două pietre și se
ascunde sub \VUgras{tufișurile ghimpoase} \textsuperscript{(21)} și sub
\VUgras{mărăcinii de păducel} \textsuperscript{(22)}.

%%K t08-tit-3 sub
\VUsaut{5.0mm}
\VUpk{10.2pt}{40}{Secerișul.}
\VUsaut{3.4mm}

%%K t08-c1-08-1 p
8. --- Pentru copiii satului strânsul grâului este o vreme foarte veselă. Se dau
vesel \VUgras{de-a berbeleacul} \textsuperscript{(24)} pe \VUgras{clăile de paie}
\textsuperscript{(23)}, îi ajută pe \VUgras{secerători} \textsuperscript{(26)}, și
în vreme ce aceștia cosesc \VUgras{holda} \textsuperscript{(27)} cu
\VUgras{coasele} lor \textsuperscript{(25)}, bine ascuțite pe \VUgras{cute}
\textsuperscript{(30)}, ei strâng grâul și îl leagă în \VUgras{snopi}
\textsuperscript{(31)} sau în \VUgras{legături} \textsuperscript{(32)}.

%%K t08-c1-09-1 p
9. --- Uneori celor mai mari li se dă să taie cu \VUgras{secera}
\textsuperscript{(12)} \VUgras{tulpinile de aur} \textsuperscript{(28)}, care poartă
\VUgras{spice} grele \textsuperscript{(29)}, pline de boabe; iar cei mici, păzind
\VUgras{vitele} \textsuperscript{(33)} care pasc pe \VUgras{livadă}
\textsuperscript{(34)} la umbra \VUgras{teiului} mare \textsuperscript{(35)},
prind cu \VUgras{fileul} lor \textsuperscript{(36)} \VUgras{fluturi} frumoși.

%%K t08-c1-09-2 p
Unii se urcă chiar în \VUgras{car} \textsuperscript{(37)}, ca să așeze snopii care
li se dau cu \VUgras{furca} \textsuperscript{(38)}.

%%K t08-c1-10-1 p
10. --- Pe toată \VUgras{câmpia} \textsuperscript{(39)}, unde își iau zborul
\VUgras{ciocârliile} \textsuperscript{(41)}, este mare forfotă. Toți sunt prinși cu
secerișul. Dar în vreme ce săracii se apucă tot de uneltele vechi ---
\VUgras{coase, seceri} și \VUgras{îmblăcie} \textsuperscript{(40)} --- proprietarii
bogați își seceră grâul sau \VUgras{secara} \textsuperscript{(43)} cu
\VUgras{secerătoarea} \textsuperscript{(42)}. Apoi, fără să ducă snopii în șură,
fac acolo pe loc \VUgras{stoguri} mari \textsuperscript{(44)}.

%%K t08-c1-11-1 p
11. --- Pe urmă se aduce \VUgras{batoza} \textsuperscript{(47)}, pe care o
învârtește \VUgras{locomobila} \textsuperscript{(45)} prin \VUgras{cureaua de
transmisie} \textsuperscript{(46)}, și astfel bobul se desparte de \VUgras{paie}
fără să părăsească ogorul unde a fost semănat.

%%K t08-tit-4 sub
\VUsaut{5.0mm}
\VUpk{10.2pt}{40}{Furtuna.}
\VUsaut{3.4mm}

%%K t08-c1-12-1 p
12. --- În vremea secerișului izbucnesc adesea \VUgras{furtuni} puternice
\textsuperscript{(53)}. Mai întâi se aude din depărtare bubuitul \VUgras{tunetului};
se ridică un \VUgras{vânt de apus} \textsuperscript{(54)} puternic și îndoaie până
la geamăt cei mai mari copaci ai \VUgras{pădurii} \textsuperscript{(49)}.
\VUgras{Nori} negri \textsuperscript{(56)} iau cerul cu asalt; îi taie zigzagurile
orbitoare ale \VUgras{fulgerului} \textsuperscript{(55)}. Începe să cadă
\VUgras{ploaia}, iar uneori și \VUgras{grindina}, culcând grâul gata de seceriș.

%%K t08-c1-13-1 p
13. --- Adesea fulgerul aprinde vreo clădire. Astfel, anul trecut acoperișul
\VUgras{morii de vânt} \textsuperscript{(50)} a ars din temelii, iar cele două
\VUgras{aripi} ale ei \textsuperscript{(51)} au fost sfărâmate.

%%K t08-c1-14-1 p
14. --- După câteva ceasuri se întoarce liniștea. La zare se ivește un
\VUgras{curcubeu} strălucitor \textsuperscript{(52)}, și firea se apucă din nou de
viața întreruptă o vreme. Sătenii, adăpostiți în \VUgras{colibe}
\textsuperscript{(48)}, se întorc la lucru și se apucă cu vitejie să dreagă paguba
abia suferită.

%%K t08-tit-5 sub
\VUsaut{6.0mm}
\VUcentre{13.2pt}{90}{\textit{A doua scenă.}}
\VUsaut{3.6mm}

%%K t08-tit-6 sub
\VUpk{10.2pt}{40}{Vânătoarea. --- Culesul viei.}
\VUsaut{1.4mm}
\VUpk{10.2pt}{40}{Pescuitul.}
\VUsaut{4.0mm}

%%K t08-c2-01-1 p
1. --- Spre sfârșitul lui \VUgras{august} sau pe la mijlocul lui
\VUgras{septembrie}, după ținut, se deschide vânătoarea. De îndată ce cele dintâi
raze de soare scot \VUgras{șopârlele} \textsuperscript{(2)} din găuri,
\VUgras{vânătorul} \textsuperscript{(5)} și iese cu \VUgras{câinii} săi
\textsuperscript{(4)}.

%%K t08-c2-02-1 p
2. --- Și-a pus \VUgras{costumul de vânătoare} trainic \textsuperscript{(9)} din
postav gros și \VUgras{jambierele} de piele, care îi vor îngădui să meargă prin
iarba udă, unde se ascund \VUgras{broaștele râioase} \textsuperscript{(1)}; și-a
luat \VUgras{pușca} \textsuperscript{(6)} și \VUgras{tolba}
\textsuperscript{(10)} --- și la drum.

%%K t08-c2-03-1 p
3. --- Focul vânătorii îl duce chiar în mijlocul viei, unde culesul este în toi.
Copiii viticultorului, care de obicei păzesc caprele lângă drum, le lasă astăzi să
pască unde vor și, legând de un par pe bătrânul \VUgras{țap} \textsuperscript{(3)},
care de bună seamă s-ar folosi de slobozenie ca să fugă, își iau și ei partea de
veselie. Unul scoate un ciorchine din \VUgras{coșul de struguri}
\textsuperscript{(12)} și se joacă stoarcând \VUgras{zeama} \textsuperscript{(14)}
într-o \VUgras{ceașcă} \textsuperscript{(13)}, ca să iasă must (vin nefermentat).
Altul se încununează cu o \VUgras{cunună de frunze} \textsuperscript{(15)}, pe care
tocmai a împletit-o. Lângă ei \VUgras{vrăbii} obraznice \textsuperscript{(16)} vin
până la teasc, ca să fure struguri.

%%K t08-c2-04-1 p
4. --- În vreme ce copiii se joacă, bărbații lucrează. \VUgras{Culegătorii}
\textsuperscript{(18)} duc strugurii în pivniță în \VUgras{coșuri de spate}
\textsuperscript{(17)}; \VUgras{dogarul} \textsuperscript{(19)} drege
\VUgras{butoaiele} \textsuperscript{(20)}, cercetează dacă \VUgras{doagele}
\textsuperscript{(22)} și \VUgras{fundurile} \textsuperscript{(21)} sunt întregi și
schimbă \VUgras{cercurile} plesnite \textsuperscript{(23)}. Tot el a făcut
\VUgras{căzile} mari \textsuperscript{(25)}, în care se varsă \VUgras{strugurii}
\textsuperscript{(26)}, ca să fermenteze.

%%K t08-c2-05-1 p
5. --- Când fermentarea se sfârșește, vinul se trage, iar rămășița se pune sub
\VUgras{teasc} \textsuperscript{(28)}; strângând încet \VUgras{șurubul} mare
\textsuperscript{(29)}, se stoarce zeama, care curge în \VUgras{cadă}
\textsuperscript{(32)} și mai dă \VUgras{vin} \textsuperscript{(31)}.

%%K t08-tit-8 sub
\VUsaut{6.0mm}
\VUpk{10.2pt}{40}{Berea și cidrul.}
\VUsaut{3.4mm}

%%K t08-c2-06-1 p
6. --- Pe zidul \VUgras{pivniței} \textsuperscript{(27)} se cațără o ramură de
\VUgras{hamei} \textsuperscript{(30)}. Aici este numai podoabă, dar în alte țări,
unde vița este foarte rară, hameiul se cultivă ca să se facă \VUgras{bere}.

%%K t08-c2-07-1 p
7. --- Micul \VUgras{măr} \textsuperscript{(33)} se pleacă sub mere, care vor da
un \VUgras{cidru} ales când se vor coace. În \VUgras{colivia} lor
\textsuperscript{(34)} \VUgras{canarul} \textsuperscript{(35)} și
\VUgras{sticletele} \textsuperscript{(36)} privesc cu pizmă vrăbiile care zburdă în
voie.

%%K t08-c2-08-1 p
8. --- Cât cuprinde ochiul, pe coasta dealurilor stau șiruri de \VUgras{araci}
\textsuperscript{(40)}, care sprijină \VUgras{vițele} minunate
\textsuperscript{(39)}, grele de \VUgras{ciorchini}.

%%K t08-c2-08-2 p
Tot anul \VUgras{viticultorul} \textsuperscript{(42)} a trudit ca să își îngrijească
\VUgras{via} \textsuperscript{(38)}, și iată-l răsplătit pentru munca lui cu o
recoltă bună. \VUgras{Culegătoarele} \textsuperscript{(41)} umplu căzile puse pe
car, tras de \VUgras{catâri} \textsuperscript{(43)}, și toți sunt veseli.

%%K t08-tit-9 sub
\VUsaut{5.0mm}
\VUpk{10.2pt}{40}{Pescuitul.}
\VUsaut{3.4mm}

%%K t08-c2-09-1 p
9. --- Tot așa și la \VUgras{pescarii cu undița} \textsuperscript{(44)}, care au
venit în zori să se așeze lângă apă cu \VUgras{undițele} lor
\textsuperscript{(45)}.

%%K t08-c2-09-2 p
\VUgras{Peștele} \textsuperscript{(47)} prinde cârligul de îndată ce lasă în apă
\VUgras{firul} \textsuperscript{(46)}, și abia apucă să schimbe \VUgras{momeala},
că o nouă jertfă se zbate deja la capătul firului.

%%K t08-c2-09-3 p
Totuși \VUgras{pescarul cu năvodul} \textsuperscript{(48)}, care își aruncă plasa
din \VUgras{barcă} \textsuperscript{(49)} spre mijlocul \VUgras{râului}
\textsuperscript{(50)}, prinde mult mai mult.

%%K t08-c2-10-1 p
10. --- Toamna este vremea plimbărilor frumoase: pe jos, \VUgras{călare}
\textsuperscript{(52)}, cu \VUgras{bicicleta} \textsuperscript{(53)}, cu
\VUgras{automobilul} \textsuperscript{(54)}. \VUgras{Drumurile}
\textsuperscript{(51)} sunt curate, și lumea se folosește de ultimele zile
frumoase, fiindcă \VUgras{rândunelele} \textsuperscript{(56)} se strâng deja pe
acoperișuri, ca să plece în țările calde. Frunzele \VUgras{alunului} bătrân
\textsuperscript{(57)} se îngălbenesc, \VUgras{alunele} cad, iar \VUgras{copacii}
înalți, care stau în \VUgras{crâng} \textsuperscript{(58)} pe \VUgras{înălțime}
\textsuperscript{(59)}, se vor despuia curând.

%%K t08-c2-11-1 p
11. --- \VUgras{Soarele} \textsuperscript{(60)} răsare mai târziu, zilele se fac mai
scurte, \VUgras{dimineața} începe să se simtă un frig destul de tăios, și iată că
stoluri de \VUgras{sitari} \textsuperscript{(61)} părăsesc ținutul nostru
neprimitor.
\VUsaut{6.0mm}
\VUfilet{22mm}
